Our system consists of two novel components that enable the acceleration of 
memcached get requests: a traffic manager and the key-value store accelerator, 
both authored in Chisel, a hardware construction language developed at Berkeley \cite{chisel}.
These components are attached to a RISC-V Rocket Core with DMA-based networking 
support. The traffic manager directs incoming network traffic to the correct 
component of our system. Incoming UDP packets that match the memcached binary 
protocol are taken out of the receive stream and handed to the key-value store 
accelerator for processing. All other packets are handed off to the DMA 
engine attached to Rocket for standard software handling.

\begin{figure}[t]
\begin{center}
\includegraphics[width=\linewidth]{../../img/system-kvstore.pdf}
\caption{Full System Design}
\label{fig:full-sys}
\end{center}
\end{figure}

When a memcached packet is handed to the accelerator, the accelerator checks 
for the presence of the requested key-value pair in its SRAM. If the pair is
available from the accelerator, a memcached binary protocol response is 
constructed and transmitted back to the client machine without involving the 
application processor. If the pair is not found in the accelerator, the packet 
is returned to the traffic manager for transmission to software on the Rocket 
Core. Once a memcached request is handed to Rocket for processing, it is 
handled as on any other system running memcached.

The cache replacement policy used by our system to determine the set of keys
placed on the acclerator is implemented entirely in software and thus is 
highly tunable without making significant modifications to either the traffic
manager or accelerator. Additionally, the accelerator is sufficiently 
generalized as to be used with any system that may benefit from handling some
requests directly at the network interface card. On the other hand, 
generalizing the traffic manager is left as future work. One avenue is to 
replace the traffic manager with a fully-programmable I/O coprocessor, allowing
on-the-fly selection of packet filtering policies.



